\section{多轮交互与宪法设计}

Compound AI 的第二步是多轮对话式数据收集。当用户只给出粗略意图时，系统需要询问细节、保持可信，并在有限轮次内给出结构化指令。

\subsection{APAC Constitution}

他们把这一需求抽象成 APAC Constitution：Accuracy、Proactivity、Adaptivity、Credibility 四个维度，每个维度都用明确的行为去量化，比如提问是否聚焦真正的关键信息、是否避免多次无关追问。
\begin{figure}[H]
\centering
\includegraphics[width=\textwidth]{frame-apac-constitution.jpg}
\caption{APAC Constitution 以多维度回应用户，实现主动询问与高信任。 \protect\footnotemark}
\end{figure}
\footnotetext{视频画面时间区间：00:16:08--00:16:22。}
\begin{importantbox}{APAC 宪法的教训}
1）Accuracy：细节必须落在用户所需的关键字段；2）Proactivity：主动发起有价值的问题；3）Adaptivity：针对不同 traveler persona 调整风格；4）Credibility：控制 hallucination，避免显眼自相矛盾。APAC 让这些目标同时优化成为可能。
\end{importantbox}
\textit{``Agent Constitution makes it possible to improve the agent performance along different axes simultaneously with very simple techniques.''}

真实对话示例是：travelers 说 ``I want to go to Hawaii'' 时，agent 先确认预算、偏好（money vs experience）、是否希望多城市，接着询问时间窗口与可接受的交通方式，并在每一轮将收集到的信息更新进 Json schema。这种逐步提问与即时验证的 pipeline 由 APAC 的四维目标驱动，保证问句既有价值又不偏离用户初始意图。

\subsection{主动问答与评估}

他们模拟了 50 个 traveler persona，每个拥有不同偏好和 hidden spreadsheet，agent 通过多轮对话补全 Json 输入，然后用 DPO fine-tune 以优化与 ground truth 的对齐。实验还将 agent 自身当作 judge，不再依赖人工全程打分，从而提高数据效率。
\begin{knowledgebox}{多轮问答指标}
关键评估包括：1）overall recall（是否收集所有信息）；2）critical recall（对最重要字段的正确性）；3）agentic score（回应的主动性与多轮连贯性）；4）DPO reward（与 judge 的 alignment）。
\end{knowledgebox}
\begin{warningbox}{避免低级人设陷阱}
主动提问必须服从目标：过度提问会浪费轮数，过度固定模板容易忽略 persona 的特殊性。模型应在 accuracy 与 efficiency 之间找到平衡。
\end{warningbox}
\begin{knowledgebox}{主动提问的启发式}
1）追踪缺失信息：每轮按字段重要度排序候选问题；2）衡量 cost-benefit：若一问带来高价值但代价高，可拆成两个小问；3）用 persona 聚焦差异：预算敏感型快速询问成本，体验导向型强调活动；4）限制轮数：超出 4 轮就考虑总结已知信息。
\end{knowledgebox}

\subsection{本章小结}

APAC Constitution 让多轮问答具备可控性和可度量性，agent judge + DPO 形成闭环，使得自监督 fine-tuning 成为可行的训练方案。

